% Generated by tools/edn_latex.py from reports/edn.md.
% Do not edit: change reports/edn.md and run `make edn`.
\ednentry{
  id = {EDN-35},
  title = {The \texorpdfstring{\texttt{download}}{download} stage owns \texorpdfstring{\texttt{data/\allowbreak{}raw/\allowbreak{}}}{data/raw/}, and the published CSV is a local cache outside the DAG},
  date = {2026-09-30},
  milestone = {M2: Reproducibility},
  topic = {Data Versioning, Pipeline Design},
  participants = {@lukas2510. Open to revisit at the next sprint review, like EDN-25.},
  decision = {The \texttt{download} stage owns \texttt{data/\allowbreak{}raw/\allowbreak{}}.
Under \texttt{download.\allowbreak{}source: zenodo} it fetches the pinned file from the Zenodo record itself and writes \texttt{data/\allowbreak{}raw/\allowbreak{}listings.\allowbreak{}parquet} as its only output, so acquisition is inside the pipeline instead of a manual step beside it.
The published CSV is kept as a local cache under \texttt{data/\allowbreak{}external/\allowbreak{}} and is neither a stage \texttt{out} nor a \texttt{dep}; \texttt{download.\allowbreak{}md5} in \texttt{params.\allowbreak{}yaml} is the dependency that pins the bytes, and the stage refuses to read a file that hashes to anything else.
The standalone pointer \texttt{data/\allowbreak{}raw/\allowbreak{}autoscout24\_\allowbreak{}dataset\_\allowbreak{}20251108.\allowbreak{}csv.\allowbreak{}dvc} is deleted, which EDN-36 records as the amendment to EDN-25.},
  rationale = {\emph{Alternatives considered:}
\begin{itemize}[leftmargin=*, topsep=2pt]
\item \textbf{Option A: the pointer stays and the stage reads it (issue \#33's option A).} \texttt{download} takes the \texttt{dvc add}-tracked CSV as a \texttt{dep} and writes only the Parquet.
\newline Pros: a fresh setup gets the CSV from DagsHub rather than from Zenodo, so it is fast and independent of Zenodo's uptime; one place to pull all data from; the snapshot's hash sits in a file at the tip of \texttt{main}, so recovering it never involves the Git history.
\newline Cons: acquisition stays a manual \texttt{dvc add} outside the pipeline, so ``reproduce from a clean clone'' is not true of the pipeline alone, which is what the data-versioning page and the course demo describe; we keep re-publishing the CSV's PII columns on a public remote; and DVC hashes every \texttt{dep} to answer \texttt{dvc status}, so 548.6 MB is read before the pipeline can say whether anything changed.
\item \textbf{Option B (chosen): the stage fetches the file, and the CSV is a local cache outside the graph.}
\newline Pros: the pipeline acquires its own input, so a clean clone reproduces it with \texttt{dvc repro} and no manual step; nothing about the CSV is re-published by us; and \texttt{dvc status} stays instant, because the stage's deps are code and params only. A clean clone does not even pay the download, because \texttt{data/\allowbreak{}raw/\allowbreak{}listings.\allowbreak{}parquet} is a pushed stage output and the stage is up to date after \texttt{dvc pull}.
\newline Cons: the file the stage really reads is not in the DAG, so DVC cannot report anything about it and a reader has to know that \texttt{download.\allowbreak{}md5} is standing in for it; and the repository stops holding a tracked copy of the raw snapshot (see the consequence below).
\item \textbf{Option C: declare the CSV as a plain \texttt{out} of the stage.}
\newline Pros: the acquired file would be in the graph, cached and recoverable with \texttt{dvc checkout}.
\newline Cons: DVC removes a stage's outputs before running it, so every \texttt{dvc repro download} would re-download 548.6 MB; and a cached \texttt{out} puts the raw PII on our remote again, which is the copy EDN-36 exists to stop.
\item \textbf{Option D: declare the CSV as a \texttt{dep}.}
\newline Cons: the same 548.6 MB hash per \texttt{dvc status} as option A, plus a worse failure mode: a missing \texttt{dep} is an error DVC raises before the stage runs, so the stage could never fetch the file it depends on, and a clean clone would be stuck.
\item \textbf{Option E: declare it as \texttt{persist: true, cache: false} (\texttt{dvc stage add -\allowbreak{}-\allowbreak{}outs-\allowbreak{}persist-\allowbreak{}no-\allowbreak{}cache}, verified present in DVC 3.67.1).}
\newline Pros: this is the option the first write-up of this decision missed, and it removes the objection to option C outright. \texttt{persist: true} means DVC does not delete the file before the run, so it is fetched once and reused, and \texttt{cache: false} keeps it out of the cache and off the remote, so there is no second copy of the PII. It would also make \texttt{dvc status} honest about the CSV: the file the stage reads would be named in the graph, and a damaged or swapped local copy would be reported by DVC and not only by our own check.
\newline Cons: a non-cached out is still hashed to build the lock and to answer \texttt{dvc status}, so it costs exactly the 548.6 MB read per status call that ruled out options A and D, on every run including the many that never touch the raw layer. It is also not recoverable despite being in the graph: \texttt{dvc checkout} cannot restore a file DVC never cached, so the lock would carry a hash with nothing behind it. And it would put a second, DVC-owned copy of a fact \texttt{download.\allowbreak{}md5} already states, so re-pinning the file would mean editing the parameter and refreshing the lock instead of only the parameter.
\end{itemize}
\par
\emph{Why:} The honest reason the CSV stays out of the graph is not that DVC cannot express ``a file the stage fetches and keeps''.
It can, with \texttt{persist: true} and \texttt{cache: false} (option E), and the first version of this entry claimed a limitation that does not exist.
The reason is what that expression costs against what it buys.
\par
What it buys is a \texttt{dvc status} that mentions the CSV and reports a local copy that has changed.
The stage already has that check and gives it more context: it re-hashes the file on every run, not only after a download, and it fails with a message naming the two things that can have happened, a damaged local copy or a genuinely re-pinned upstream file, because those need different answers.
And because \texttt{download.\allowbreak{}md5} is a parameter, re-pinning is a change DVC reruns on, which is the dependency edge a graph entry would have drawn.
\par
What it costs is a 548.6 MB read on every \texttt{dvc status}, which is the cost that ruled out the \texttt{dep} shape in the first place, plus a lock entry with no recoverable object behind it.
So the property is already covered by the stage while the cost would be paid by everyone on every status call, and the file stays a cache.
\par
It lives under \texttt{data/\allowbreak{}external/\allowbreak{}} because that is the third-party slot of the project layout: a copy of a published artefact, reproducible from the DOI and the pinned MD5, with nothing for DVC to version. Keeping it out of \texttt{data/\allowbreak{}raw/\allowbreak{}} is what lets that directory belong to the pipeline alone.
\par
\textbf{Consequence, recorded rather than argued away.} After this change no \texttt{.\allowbreak{}dvc} file at the tip of the branch references the raw blob.
The blob is still on DagsHub and in local caches, and the pointer's hash is still in the Git history of the deleted file, so an earlier commit can still \texttt{dvc pull} and the snapshot is recoverable.
But recovering it without Zenodo now means somebody reading a hash out of the Git history, and a workspace-scoped \texttt{dvc gc} would delete the blob without asking.
That is why the rule never to run \texttt{dvc gc} without \texttt{-\allowbreak{}-\allowbreak{}all-\allowbreak{}commits} is written down in the data-versioning page as part of this change.},
  ai = {Information seeking, Alternative generation, Alternative assessment, Recommendation, Solution generation},
  response = {Accepted with modifications},
  assessment = {AI implemented issue \#33's option B and proposed the refinement that the CSV is a local cache rather than a stage output, with the measurements behind it: the CSV is 548.6 MB, the Parquet 215.3 MB, the stage 88 s cold and 66 s warm.
Its option table was wrong in one way that mattered: it argued the CSV could only be an \texttt{out} that is re-downloaded every run or a \texttt{dep} that is hashed on every \texttt{dvc status}, and presented the cache as the only way out, which made a DVC limitation out of a design trade-off.
An adversarial review of PR \#54, also run with AI, found \texttt{-\allowbreak{}-\allowbreak{}outs-\allowbreak{}persist-\allowbreak{}no-\allowbreak{}cache} in DVC 3.67.1 and established that a \texttt{persist: true, cache: false} out is neither deleted before the run nor pushed, and the same review found the recoverability consequence above.
Lukas kept the mechanism and required the rationale to be rewritten so that it rests on the cost of the 548.6 MB hash and on the check the stage already performs, rather than on a limitation that does not exist, and required the consequence to be recorded.},
  aievidence = {Claude Code session on 2026-09-30 implementing \href{https://github.com/mlops-2627q1-mds-upc/MLOPS_Recommenditos/issues/33}{issue \#33}, which produced the first option table in the body of \href{https://github.com/mlops-2627q1-mds-upc/MLOPS_Recommenditos/pull/54}{PR \#54}; a second Claude Code session the same day reviewing that PR adversarially, which produced the \texttt{persist} finding and the recoverability consequence; Lukas directed that the entry be corrected rather than the mechanism changed.},
  evidence = {\href{https://github.com/mlops-2627q1-mds-upc/MLOPS_Recommenditos/issues/33}{Issue \#33}; \href{https://github.com/mlops-2627q1-mds-upc/MLOPS_Recommenditos/pull/54}{PR \#54}; the \texttt{download} stage in \texttt{dvc.\allowbreak{}yaml} and the \texttt{download} block of \texttt{params.\allowbreak{}yaml}; \href{https://github.com/mlops-2627q1-mds-upc/MLOPS_Recommenditos/blob/main/docs/docs/data-versioning.md}{Data versioning}, ``Raw data''; \texttt{tests/\allowbreak{}test\_\allowbreak{}download.\allowbreak{}py}; EDN-25, EDN-34, EDN-36.},
}
